\documentclass[11pt]{article}
\usepackage{shared/template_tgir_article}
\usepackage{booktabs}
\usepackage{longtable}
\usepackage{tabularx}
\usepackage{array}
\usepackage{enumitem}
\usepackage{xcolor}
\usepackage{listings}

\lstset{
    basicstyle=\ttfamily\small,
    breaklines=true,
    columns=fullflexible,
    frame=single
}

\newcommand{\product}{TGIR QuantLib Tools}
\newcommand{\code}[1]{\texttt{#1}}
\newcommand{\docdate}{2026-04-06}

% Visual design is controlled exclusively by template_tgir_article.sty and
% IEEEtran. Keep this file limited to content-support packages and commands.


\title{Developer Guide\\\large \product}
\author{TG Investments and Research}
\date{\docdate}

\begin{document}
\maketitle
\tableofcontents
\newpage

\section{Audience And Goal}
This guide is for developers who need to understand the application structure, local workflow, extension points, and delivery expectations.

\section{Repository Layout}
\begin{longtable}{p{0.28\textwidth}p{0.66\textwidth}}
\toprule
Path & Responsibility \\
\midrule
\path{app.py} & Local Flask entrypoint. \\
\path{portfolio.py} & QuantLib market construction, trade pricing, and analytics logic. \\
\path{tgir_quantlib_tools/app_factory.py} & Flask app construction. \\
\path{tgir_quantlib_tools/routes.py} & Route registration and HTTP handlers. \\
\path{tgir_quantlib_tools/dashboard.py} & Session state, dashboard payload shaping, and trade-editor context building. \\
\path{tgir_quantlib_tools/quantlib_model.py} & Data-model page and research section context. \\
\path{tgir_quantlib_tools/research.py} & Curated paper metadata for swaption and cliquet references. \\
\path{templates/} & Jinja templates for the UI. \\
\path{tests/} & Smoke, calibration, and identity tests. \\
\path{docs/latex/} & Audience-specific LaTeX documentation. \\
\bottomrule
\end{longtable}

\section{Application Flow}
\begin{enumerate}[leftmargin=1.5em]
    \item \code{create\_app()} loads configuration and builds the Flask app.
    \item Authentication and template filters are registered.
    \item Requests hit \path{tgir_quantlib_tools/routes.py}.
    \item Protected routes load session state from \path{tgir_quantlib_tools/dashboard.py}.
    \item Pricing functions in \path{portfolio.py} normalize state and construct QuantLib objects.
    \item Templates render the final dashboard, editor, or model/research page.
\end{enumerate}

\section{State Model}
The session state is a nested dictionary with:
\begin{itemize}[leftmargin=1.5em]
    \item \code{market} for rates, swaption volatility, Bermudan seed, and SPX equity assumptions,
    \item \code{trades.swap},
    \item \code{trades.european\_swaption},
    \item \code{trades.bermudan\_swaption}, and
    \item \code{trades.equity\_cliquet}.
\end{itemize}
Always pass data through \code{normalize\_portfolio\_state()} before constructing pricing objects.

\section{Pricing Responsibilities}
\subsection{portfolio.py}
This is the core quantitative module. It owns:
\begin{itemize}[leftmargin=1.5em]
    \item curve normalization and construction,
    \item swaption-volatility matrix normalization,
    \item swap, swaption, and cliquet instrument creation,
    \item Bermudan Gaussian short-rate calibration,
    \item price tables,
    \item risk ladders, and
    \item cliquet scenario and Monte Carlo analytics.
\end{itemize}

\subsection{dashboard.py}
This module should stay presentation-oriented. It shapes payloads, form definitions, and display-specific summaries, but should not duplicate pricing logic already available in \path{portfolio.py}.

\section{How To Add Or Change A Trade}
\begin{enumerate}[leftmargin=1.5em]
    \item Extend \code{default\_portfolio\_state()} and normalization rules.
    \item Add a human-readable entry in \code{TRADE\_SEQUENCE} and \code{TRADE\_TITLES}.
    \item Implement instrument creation and pricing in \path{portfolio.py}.
    \item Add form definitions and update handlers in \path{tgir_quantlib_tools/dashboard.py}.
    \item Add template support in \path{templates/dashboard.html} and \path{templates/trade\_form.html}.
    \item Extend smoke and regression tests before merging.
\end{enumerate}

\section{Local Development Workflow}
\subsection{Setup}
\begin{lstlisting}
python3 -m venv .venv
./.venv/bin/python -m pip install -r requirements.txt
cp .env.example .env
\end{lstlisting}

\subsection{Useful Commands}
\begin{lstlisting}
./.venv/bin/python app.py
./.venv/bin/python build_SOFR_curve.py
./.venv/bin/python -m unittest discover -s tests
./.venv/bin/python -m py_compile portfolio.py tgir_quantlib_tools/*.py
\end{lstlisting}

\section{Documentation Workflow}
The LaTeX docs are built with:
\begin{lstlisting}
make -C docs/latex
\end{lstlisting}
Keep audience-specific material separated:
\begin{itemize}[leftmargin=1.5em]
    \item end users: workflows and UI,
    \item quants: models and assumptions,
    \item developers: architecture and extension points,
    \item IT: runtime operations,
    \item testing: regression methodology,
    \item deployment: CI/CD and DigitalOcean procedures.
\end{itemize}

\section{Testing Expectations}
Changes that touch pricing must update tests in one of three layers:
\begin{itemize}[leftmargin=1.5em]
    \item \path{tests/test\_sofr\_curve.py} for calibration or rates logic,
    \item \path{tests/test\_cliquet.py} for cliquet payoff identities, and
    \item \path{tests/test\_smoke.py} for route and rendering behavior.
\end{itemize}

\section{Change Review Checklist}
\begin{itemize}[leftmargin=1.5em]
    \item Did the change preserve state normalization?
    \item Is new pricing logic in \path{portfolio.py} rather than duplicated in templates or routes?
    \item Do calibration instruments still reprice?
    \item Are tests deterministic?
    \item Are docs updated in both Markdown and LaTeX where relevant?
\end{itemize}

\end{document}
